\documentclass[12pt,letterpaper]{article}
% Paquetes esenciales
\usepackage[utf8]{inputenc}
\usepackage[spanish]{babel}
\usepackage[margin=2.5cm]{geometry}
\usepackage{hyperref}
\usepackage{xcolor}
\usepackage{listings}
% Configuración de colores para bloques de código
\definecolor{codegray}{rgb}{0.95,0.95,0.95}
\definecolor{codeblue}{rgb}{0.1,0.2,0.8}
\definecolor{codegreen}{rgb}{0.1,0.5,0.1}
\lstset{
backgroundcolor=\color{codegray},
basicstyle=\ttfamily\footnotesize,
keywordstyle=\color{codeblue}\bfseries,
stringstyle=\color{codegreen},
frame=single,
breaklines=true,
breakatwhitespace=false,
captionpos=b,
showstringspaces=false
}
\title{\textbf{Guía de Despliegue \\ Sistema de Control de Acceso Autónomo}}
\author{Saúl Quesada Soto}
\date{Octubre 2026}
\begin{document}
\maketitle
\section*{Resumen del Sistema}
Este documento detalla el despliegue de la capa de Yocto Project (\texttt{meta-control-acceso}) que opera un sistema de control de acceso físico operado por Inteligencia Artificial. Diseñado para ejecutarse en una \textbf{Raspberry Pi 4} de forma autónoma (Zero-Touch Boot).
El sistema utiliza \textbf{GStreamer} para transmisión de video concurrente de baja latencia y \textbf{OpenCV DNN} (YuNet / MobileNet-SSD) para detección facial y prevención de \textit{Tailgating}.
\section{Requisitos de Hardware}
\begin{itemize}
\item \textbf{Placa:} Raspberry Pi 4 Model B.
\item \textbf{Cámara:} Cámara Web USB Genérica (Ej. Fobos Redragon).
\item \textbf{Actuadores GPIO:}
\begin{itemize}
\item LED Rojo (Bloqueo) 
→
→
 Pin físico 15 (GPIO 22).
\item LED Verde (Apertura) 
→
→
 Pin físico 13 (GPIO 27).
\item Botón REX (Salida manual) 
→
→
 Pin físico 11 (GPIO 17) conectado a Tierra (GND).
\end{itemize}
\item \textbf{Memoria:} MicroSD de al menos 16GB.
\end{itemize}
\section{Guía de Despliegue y Compilación}
\subsection{Preparación del Entorno Host (Ubuntu)}
El sistema debe compilarse en una distribución Linux (preferiblemente Ubuntu 22.04 o 24.04). Instale las dependencias base de Yocto y la herramienta de flasheo rápido \texttt{bmaptool}:
\begin{lstlisting}[language=bash]
sudo apt update
sudo apt install gawk wget git diffstat unzip texinfo gcc build-essential chrpath socat cpio python3 python3-pip python3-pexpect xz-utils debianutils iputils-ping python3-git python3-jinja2 libegl1-mesa libsdl1.2-dev pylint3 xterm python3-subunit mesa-common-dev zstd liblz4-tool bmap-tools
\end{lstlisting}
\subsection{Inicializar el Entorno Yocto}
Asegúrese de haber descargado Poky y las capas necesarias (\texttt{meta-raspberrypi}, \texttt{meta-openembedded}). Clone esta capa en su directorio \texttt{layers/}.
\begin{lstlisting}[language=bash]
Inicializar el entorno de construccion de Yocto
source layers/poky/oe-init-build-env build/
\end{lstlisting}
\subsection{Configuración Dinámica (Variables de Entorno)}
Antes de compilar, configure la IP de la computadora del vigilante (quien recibirá el video) y el tipo de codificador. Estos valores pueden modificarse posteriormente editando el archivo \texttt{/etc/default/control-acceso} vía SSH sin necesidad de recompilar el sistema.
Ruta del archivo fuente: \texttt{recipes-acceso/control-acceso/control-acceso/control-acceso.env}
\begin{lstlisting}[language=bash]
Usar "produccion" para activar el codificador de hardware H.264
ACCESO_FUENTE=produccion
ACCESO_CAM=/dev/video0
IP de la computadora Windows que recibira el video
ACCESO_HOST=172.21.255.122
ACCESO_GPIOCHIP=/dev/gpiochip0
\end{lstlisting}
\subsection{Compilación}
Ejecute la construcción de la imagen base. Yocto instalará Python, GStreamer, OpenCV, inyectará los modelos de IA y configurará el demonio \texttt{systemd}.
\begin{lstlisting}[language=bash]
bitbake core-image-base
\end{lstlisting}
\subsection{Flasheo de la Memoria SD}
Inserte la MicroSD en su computadora. Identifique la ruta de su memoria (ej. \texttt{/dev/sda}) usando \texttt{lsblk}. Desmonte cualquier partición automática y flashee el sistema usando \texttt{bmaptool} (superior a \texttt{dd} en velocidad y seguridad):
\begin{lstlisting}[language=bash]
sudo umount /dev/sda* 2>/dev/null ; sudo bmaptool copy tmp/deploy/images/raspberrypi4-64/core-image-base-raspberrypi4-64.rootfs.wic.bz2 /dev/sda
\end{lstlisting}
\section{Operación del Sistema (Zero-Touch)}
\begin{enumerate}
\item Inserte la tarjeta SD en la Raspberry Pi 4.
\item Conecte la Cámara USB y el cable Ethernet a la red local.
\item Conecte la alimentación (USB-C).
\item \textbf{No interactúe con la placa ni entre por SSH.} El sistema \texttt{systemd} esperará automáticamente 5 segundos a que la cámara despierte en el bus USB y arrancará el pipeline en modo \textit{headless}.
\end{enumerate}
\subsection{Recepción del Video (Cliente Windows/Linux)}
En la computadora del vigilante (cuya IP fue configurada previamente), abra una consola CMD/Bash y ejecute el siguiente comando para recibir el video UDP de latencia ultra baja:
\begin{lstlisting}[language=bash]
gst-launch-1.0 -v udpsrc port=5000 ! "application/x-rtp, media=video, clock-rate=90000, encoding-name=H264, payload=96" ! rtph264depay ! h264parse ! avdec_h264 ! videoconvert ! autovideosink sync=false
\end{lstlisting}
\subsection{Persistencia y Evidencia Legal}
Cualquier evento de intrusión, salida manual (\textit{Override}) o \textit{Tailgating} (múltiples personas) será registrado en la bitácora y se capturará una fotografía. Esta evidencia se guarda de forma persistente, a prueba de cortes eléctricos o reinicios, en la partición ext4 física:
\begin{itemize}
\item Fotografías: \texttt{/home/root/control-acceso/evidencias/}
\item Bitácora: \texttt{/home/root/control-acceso/registro.csv}
\end{itemize}
\end{document}
